\documentclass[11pt,a4paper]{article}
\usepackage[T1]{fontenc}
\usepackage[utf8]{inputenc}
\usepackage[brazil]{babel}
\usepackage{lmodern}
\usepackage[margin=2.25cm]{geometry}
\usepackage{amsmath,booktabs,caption,graphicx,microtype}
\usepackage[hidelinks]{hyperref}

\hypersetup{
  pdftitle={Método de Newton para o sistema tridiagonal de Broyden},
  pdfauthor={Guilherme Fontana e Pedro Stahel de Azevedo},
  pdfsubject={CI1164 — Introdução à Computação Científica}
}
\captionsetup{font=small,labelfont=bf,hypcap=false}

\title{\textbf{Método de Newton para o sistema tridiagonal de Broyden}\\[0.25em]
       \large Comparação entre eliminação densa, em banda e pelas extremidades}
\author{Guilherme Fontana (\texttt{gsf25}) \and Pedro Stahel de Azevedo (\texttt{psa25}, GRR20252084)}
\date{Departamento de Informática, Universidade Federal do Paraná\\
      CI1164 — Introdução à Computação Científica\\
      18 de novembro de 2026}

\newcommand{\painel}[3]{%
  \begin{minipage}[t]{0.485\linewidth}
    \vspace{0pt}%
    \centering
    \includegraphics[width=\linewidth]{#1}
    \captionof{figure}{#2}\label{#3}
  \end{minipage}}
\newcommand{\tempo}{resultados/integrada/graficos/tempo-}
\newcommand{\hardware}{evidencias/gsf25/resultados/graficos/}

\begin{document}
\maketitle

\begin{abstract}
Comparamos três implementações do método de Newton aplicado ao sistema
tridiagonal de Broyden. A referência usa uma matriz densa, a segunda
implementação usa eliminação em banda com pivoteamento parcial e a
terceira elimina pelas duas extremidades quando há uma garantia sobre
os pivôs. Em um Intel Core i5-7500, com 25 passos e $N=1000$, os tempos
médios do cálculo de Newton foram de $3{,}740$~s, $0{,}401$~ms e
$0{,}151$~ms, respectivamente. Em $N=20000$, o tempo caiu de
$8{,}06$~ms em banda para $3{,}03$~ms na versão otimizada. Os vetores
finais concordaram nas dimensões submetidas às mesmas verificações
numéricas. Uma campanha anterior com LIKWID complementa a análise de
cache, banda e vazão aritmética na comparação entre a referência densa
e a primeira implementação tridiagonal.
\end{abstract}

\noindent\textbf{Palavras-chave.} Método de Newton, sistema de Broyden,
matriz tridiagonal, desempenho.

\section{Introdução}

Quando a Jacobiana de um sistema não linear é tridiagonal, é possível
construí-la e resolver o sistema de Newton com trabalho proporcional ao
número de incógnitas. Essa mudança de estrutura explica a primeira parte
do ganho medido. Dentro do custo linear, o número de acessos à memória e
a dependência entre divisões ainda influenciam o tempo. Este trabalho
mede esses dois efeitos com três implementações em C que compartilham
o mesmo laço de Newton, a mesma interface e o mesmo relógio.

\section{Formulação}

Seja $x=(x_1,\ldots,x_N)^\mathsf{T}$, com vizinhos de contorno
$x_0=x_{N+1}=0$. O sistema é $F(x)=0$, onde
\begin{equation}
  F_i(x)=(3-2x_i)x_i-x_{i-1}-2x_{i+1}+c_i,\qquad 1\leq i\leq N.
  \label{eq:broyden}
\end{equation}
Usamos $c_i=1$ para $i<N$. A campanha integrada adota $c_N=0$, como no
gerador preservado com a referência do trabalho. Na formulação clássica
de Moré, Garbow e Hillstrom \cite{more}, $c_N=1$. A figura oficial do
enunciado estava indisponível para conferência. A opção
\texttt{LAST\_CONSTANT} permite compilar as duas variantes, que têm
a mesma Jacobiana. O parâmetro de entrada que inicializa todas as
incógnitas é denotado por $a$ para distingui-lo do vizinho de contorno.

A diagonal principal da Jacobiana é $d_i=3-4x_i$. A subdiagonal vale
$-1$ e a superdiagonal vale $-2$. A cada passo $k$, calculamos $s$ por
$J(x^{(k)})s=-F(x^{(k)})$ e fazemos $x^{(k+1)}=x^{(k)}+s$.
Com $\varepsilon>0$, o programa verifica se $\|F(x^{(k)})\|_\infty
\leq\varepsilon$ antes da resolução linear. Com $\varepsilon=0$, o
laço prossegue até o limite de iterações, sujeito a falha numérica.

\clearpage
\section{Implementações}

\paragraph{Referência densa.} A matriz ocupa um bloco de $N^2$ valores,
com um vetor de ponteiros para as linhas. A Jacobiana é montada por
inteiro, inclusive os zeros estruturais. A eliminação de Gauss busca
o maior pivô em módulo na coluna ativa e atualiza toda a submatriz
restante. As trocas de linha permutam ponteiros. No termo dominante,
a eliminação faz aproximadamente $2N^3/3$ operações de soma e
multiplicação.

\paragraph{Eliminação em banda.} A construção da Jacobiana calcula
a diagonal principal. A subdiagonal é conhecida e a primeira
superdiagonal recebe $-2$ no início da eliminação. O pivoteamento
compara a linha corrente com a seguinte. Uma troca pode criar uma
segunda superdiagonal, prevista no armazenamento e usada na
retrossubstituição. Essa estrutura também é usada na rotina
\texttt{DGTSV} do LAPACK \cite{lapack}.

\paragraph{Eliminação pelas extremidades.} Quando todos os $d_i$ são
finitos e $|d_i|\geq 3$, a versão otimizada usa duas recorrências
independentes até a linha central. A partir da esquerda, o pivô
ajustado é $p_i=d_i-2/p_{i-1}$. A partir da direita, a recorrência
usa $p_{i+1}$. O primeiro pivô de cada lado tem módulo de pelo menos
$3$. Se o pivô anterior tem módulo de pelo menos $2$, seu efeito no
seguinte é de no máximo $1$ e o novo pivô tem módulo de pelo menos
$2$. A linha central recebe as duas contribuições e seu pivô tem
módulo de pelo menos $1$.

O código calcula o inverso de cada pivô uma vez e guarda os
coeficientes necessários para recuperar a solução da linha central
até as extremidades. A intercalação das duas cadeias permite
sobrepor parte da latência das divisões escalares no mesmo núcleo.
Na ausência da condição de pivôs, o programa usa a eliminação em
banda com trocas. A versão otimizada também separa as bordas no
cálculo de $-F$, alinha os vetores a 64 bytes e informa esse
alinhamento ao compilador nos laços apropriados.

\begin{center}
\captionof{table}{Ordem do custo por atualização de Newton.}
\begin{tabular}{lccc}
\toprule
Implementação & Construção de $J$ & Resolução linear & Memória de trabalho \\
\midrule
Densa & $O(N^2)$ & $O(N^3)$ & $O(N^2)$ \\
Em banda & $O(N)$ & $O(N)$ & $O(N)$ \\
Otimizada & $O(N)$ & $O(N)$ & $O(N)$ \\
\bottomrule
\end{tabular}
\end{center}

\section{Protocolo experimental}

As medições integradas usaram Linux x86-64 e GCC 14.2.0, com
\texttt{-O3 -march=native -mavx -ffp-contract=fast}. O processador
Intel Core i5-7500 \cite{intel} tem L1 de dados de 32~KiB, L2 de
256~KiB por núcleo e L3 compartilhada de 6~MiB na estação medida.
As execuções foram fixadas na CPU lógica 2. Usamos $a=-1$,
$\varepsilon=0$ e 25 atualizações para as dimensões 32, 64, 128, 256,
512, 1000, 2000, 4000, 8000, 9000, 10000 e 20000.
A coleta densa foi limitada a $N\leq1000$ pelo custo cúbico.
Cada combinação teve um aquecimento descartado e dez repetições.
A ordem das dimensões foi embaralhada a cada rodada e a ordem das
versões foi alternada.

O relógio \texttt{CLOCK\_MONOTONIC\_RAW} registra o tempo acumulado de
Newton e das duas regiões internas, Jacobiana e sistema linear.
O cálculo de Newton é medido a cada atualização. A saída é
formatada entre as regiões medidas e descartada em \texttt{/dev/null}.
O tempo de parede inclui a formatação, a inicialização e o encerramento.
Os resultados de tempo usam média e desvio-padrão amostral.

\clearpage
\section{Validação e resultados de tempo}

A suíte compara os passos até $N=10$ com uma implementação densa
independente em Python e verifica o resíduo do sistema linear.
As três versões concordaram nos vetores finais dos casos em comum
até $N=512$, com tolerâncias absoluta de $2\cdot10^{-12}$ e relativa
de $2\cdot10^{-11}$. As duas tridiagonais tiveram resíduo final
máximo de $1{,}11\cdot10^{-15}$ na série de 25 passos. Também foram
verificados casos com troca de linha, matriz singular, entradas
inválidas e limite de zero atualizações. Os testes foram executados
para os dois valores de $c_N$.

\medskip
\noindent
\painel{\tempo newton_total.pdf}{Tempo acumulado de Newton nas 25 atualizações.}{fig:tempo-newton}
\hfill
\painel{\tempo jacobiana.pdf}{Tempo de construção das Jacobianas.}{fig:tempo-j}
\par\medskip
\noindent
\painel{\tempo sistema_linear.pdf}{Tempo de resolução dos sistemas lineares.}{fig:tempo-linear}
\hfill
\begin{minipage}[t]{0.47\linewidth}
  \vspace{0pt}
  \captionsetup{skip=6pt}
  \captionof{table}{Tempo de Newton em ms. Os valores são média $\pm$ desvio-padrão.}
  \label{tab:tempo}
  \centering\small
  \begin{tabular}{@{}lcc@{}}
  \toprule
  Versão & $N=1000$ & $N=20000$ \\
  \midrule
  Densa & $3740\pm420$ & \textemdash \\
  Em banda & $0{,}401\pm0{,}027$ & $8{,}06\pm0{,}22$ \\
  Otimizada & $0{,}151\pm0{,}014$ & $3{,}03\pm0{,}31$ \\
  \bottomrule
  \end{tabular}
  \par\smallskip
  \footnotesize O traço indica ausência de medição.
\end{minipage}

\medskip
Nos gráficos de tempo, os dois eixos estão em escala logarítmica.
A faixa em torno da média representa um desvio-padrão amostral.
Em $N=1000$, a mudança da matriz densa para a banda reduziu o
tempo médio de Newton em cerca de $9\,300$ vezes. A resolução linear
respondeu por $99{,}4\%$ do tempo da referência densa nesse ponto
(Figuras~\ref{fig:tempo-newton} a \ref{fig:tempo-linear}). Em
$N=20000$, a matriz densa ocuparia $3{,}2$~GB só para os
coeficientes. As versões tridiagonais reservam cinco vetores, que
somam $800$~kB nessa dimensão.

Em $N=20000$, a otimizada foi $2{,}66$ vezes mais rápida que a
versão em banda no cálculo de Newton. O tempo linear caiu de
$6{,}73$ para $1{,}97$~ms. As médias da Jacobiana foram próximas,
$0{,}335$ e $0{,}350$~ms. Em um ensaio anterior com $c_N=1$,
a mudança isolada da eliminação por uma ponta para duas reduziu
o tempo linear de $3{,}54$ para $1{,}99$~ms no mesmo tamanho.
Esse ensaio consta nos dados complementares e ajuda a atribuir
parte do ganho à organização das divisões.

No processo completo, as medianas foram $136{,}39$~ms em banda
e $131{,}15$~ms na otimizada, também em $N=20000$. Cada execução
formata $520\,000$ linhas de valores, somando o vetor inicial e as
25 atualizações. A formatação tem peso mesmo com a saída
descartada, o que limita o ganho de tempo de parede.

\clearpage
\section{Contadores de hardware}

Os contadores das Figuras~\ref{fig:l2-newton} a \ref{fig:mem-linear}
vêm de uma campanha anterior com LIKWID 5.5.2 \cite{likwid}, na
CPU lógica 3 de uma estação com o mesmo modelo de processador.
\emph{Baseline} designa a referência densa original e
\emph{otimizada} designa a primeira versão tridiagonal, com
pivoteamento e quatro vetores de coeficientes. A campanha usou
$c_N=0$, $a=-1$ e 25 passos, com uma execução por grupo, dimensão
e versão. Há pontos densos até $N=2000$ e tridiagonais até $N=20000$.
Os CSVs brutos e os fontes correspondentes acompanham o relatório.

A coleta integrada usou os relógios do programa. O valor de
\texttt{perf\_event\_paranoid} era 3 e \texttt{/dev/cpu} estava ausente.
Uma sonda anterior com LIKWID 5.4.1 no ambiente
restrito registrou bloqueio de acesso à unidade de monitoramento.
A eliminação pelas extremidades segue sem uma série de contadores.
As medidas desta seção descrevem apenas os executáveis da campanha
anterior.

\subsection{Faltas na cache L2}

\noindent
\painel{\hardware l2_newton.pdf}{Fração de solicitações a L2 que resultaram em falta no cálculo de Newton.}{fig:l2-newton}
\hfill
\painel{\hardware l2_jacobiana.pdf}{Fração de faltas a L2 na construção da Jacobiana.}{fig:l2-j}
\par\medskip
\noindent
\painel{\hardware l2_sistema_linear.pdf}{Fração de faltas a L2 na resolução linear.}{fig:l2-linear}

\medskip
Na resolução linear densa, a fração de faltas foi $0{,}0181$ em
$N=128$ e $0{,}2075$ em $N=2000$ (Figura~\ref{fig:l2-linear}).
A matriz de $128$~KiB em $N=128$ cabe na L2 de $256$~KiB. Em
$N=2000$, ela ocupa $32$~MB e ultrapassa também a capacidade de
L3. A fração tridiagonal em $N=2000$ foi $0{,}2119$ na mesma
região. Frações próximas podem corresponder a números de
solicitações e de faltas muito diferentes. Como cada ponto
vem de uma execução, diferenças pequenas merecem cautela.

\clearpage
\subsection{Operações de ponto flutuante}

\noindent
\painel{\hardware flops_newton.pdf}{Vazão aritmética em Newton.}{fig:dp-newton}
\hfill
\painel{\hardware flops_jacobiana.pdf}{Vazão aritmética na Jacobiana.}{fig:dp-j}
\par\medskip
\noindent
\painel{\hardware flops_sistema_linear.pdf}{Vazão aritmética na resolução linear.}{fig:dp-linear}

\medskip
As séries \texttt{FLOPS\_DP} e \texttt{FLOPS\_AVX} das
Figuras~\ref{fig:dp-newton} a \ref{fig:dp-linear} vieram de
execuções independentes. Para estimar a participação das
operações AVX dentro de uma execução, usamos as métricas
\emph{DP} e \emph{AVX DP} do próprio grupo \texttt{FLOPS\_DP}.
Em $N=2000$, a resolução linear densa registrou
$2684$~MFLOP/s em precisão dupla, dos quais $2676$~MFLOP/s
vieram de operações empacotadas de 256 bits. A eliminação
tridiagonal com pivoteamento registrou $537$~MFLOP/s na
mesma região e dimensão. A densa alcança vazão aritmética
maior em laços que operam sobre várias colunas consecutivas.
A eliminação em banda resolve o sistema em menos tempo porque
faz um número muito menor de operações.

Na Jacobiana tridiagonal, em $N=1000$, a taxa medida foi de
$1024$~MFLOP/s em precisão dupla e $1015$~MFLOP/s em AVX DP
dentro da mesma execução. Isso é compatível com a vetorização
do cálculo de $3-4x_i$ em um laço com acessos consecutivos.
A montagem densa passa a maior parte da região gravando
constantes e zeros, atividades que demandam tráfego de
memória e poucas operações de ponto flutuante.

\clearpage
\subsection{Banda de memória}

\noindent
\painel{\hardware mem_newton.pdf}{Banda de DRAM reportada para Newton.}{fig:mem-newton}
\hfill
\painel{\hardware mem_jacobiana.pdf}{Banda de DRAM reportada para a Jacobiana.}{fig:mem-j}
\par\medskip
\noindent
\painel{\hardware mem_sistema_linear.pdf}{Banda de DRAM reportada para a resolução linear.}{fig:mem-linear}
\hfill
\begin{minipage}[t]{0.475\linewidth}
  \vspace{0pt}\small
  O grupo \texttt{MEM} contabiliza leituras e escritas de DRAM no
  soquete. Em regiões de poucos microssegundos, tráfego de fundo e
  instrumentação podem afetar a leitura. A linha pontilhada marca
  $38{,}4$~GB/s, teto nominal de dois canais DDR4-2400 \cite{intel}.
  Marcadores vazados indicam valores acima desse teto. Em $N=32$,
  a Jacobiana tridiagonal reportou $2{,}10$~TB/s, um valor
  incompatível com a banda da plataforma.
\end{minipage}

\medskip
Na referência densa com $N=2000$, o volume registrado para
Newton foi $268{,}95$~GB e a sua região interna de resolução
linear registrou $1091{,}66$~GB. O ponto marcado com uma cruz
na Figura~\ref{fig:mem-newton} tem banda abaixo do teto e volume
inconsistente. A descrição do grupo prevê a possibilidade de
transbordamento dos contadores em intervalos longos. Esses
pontos foram preservados nos dados e excluídos da interpretação
quantitativa de ganhos em banda.

\section{Conclusão}

A estrutura tridiagonal permitiu reduzir o custo por atualização
de $O(N^3)$ para $O(N)$ na resolução linear e de $O(N^2)$ para
$O(N)$ no armazenamento. A eliminação pelas extremidades
reduziu o custo dentro do regime linear, com ganho de $2{,}66$
vezes no cálculo de Newton em $N=20000$ e concordância
numérica nos casos verificados. O ganho de tempo de parede
foi menor devido à formatação dos vetores. Os contadores da
campanha anterior ajudam a interpretar a referência densa e
a eliminação em banda. Uma medição de hardware dos binários
integrados é necessária para caracterizar a eliminação pelas
extremidades com as mesmas métricas.

\small
\begin{thebibliography}{9}
\bibitem{more}
J. J. Moré, B. S. Garbow e K. E. Hillstrom.
\textit{Testing Unconstrained Optimization Software}. ACM Transactions on
Mathematical Software, 7(1), 17–41, 1981.
\href{https://doi.org/10.1145/355934.355936}{DOI 10.1145/355934.355936}.

\bibitem{lapack}
Netlib. \textit{LAPACK: DGTSV}.
\href{https://netlib.org/lapack/explore-html/d6/ddb/group__gtsv_ga6b06a11b8e2dced10bc9cd3c20a09484.html}{Documentação da rotina}.

\bibitem{intel}
Intel. \textit{Intel Core i5-7500 Processor Specifications}.
\href{https://www.intel.com/content/www/us/en/products/sku/97123/intel-core-i57500-processor-6m-cache-up-to-3-80-ghz/specifications.html}{Ficha técnica}.

\bibitem{likwid}
J. Treibig, G. Hager e G. Wellein.
\textit{LIKWID: A Lightweight Performance-Oriented Tool Suite for x86 Multicore Environments}.
2010. \href{https://arxiv.org/abs/1004.4431}{arXiv:1004.4431}.
\end{thebibliography}
\end{document}
